% MAIN SNAPSHOT — LSPV/RMV::0948::2026-10-02T09:48:00-07:00
% !TEX program = xelatex
\documentclass[11pt]{article}
\usepackage[margin=.86in]{geometry}
\usepackage{amsmath,amssymb}
\begin{document}
\section*{LSPV-5 MAIN Snapshot 0948}

\textbf{STATUS: OPEN.} P versus NP remains open. No unrestricted SAT lower bound,
arbitrary-algorithm acquisition theorem, nonrelativizing hinge, or nonalgebrizing hinge is established.

\section*{CORPUS RECONCILIATION}
\textbf{Corpus surveyed.}
The repository tree, root \texttt{main.tex}, AGENTS/CODEX/pnp-research governance, TASK-001--007,
verified MAIN history through 0858 and the immediately preceding 0945 draft, materially relevant
LSPV-1--4 reports with predecessor LSPV-4 TRACECLOSE1 preflighted, post-0858 LSPV-2 TURANEX1,
PRIME-LQP-1, PRIMECRT1, BITRANK1, HEC/VSC-HEC/G-PPC-HEC/CRS-PPC/CCSPC/RCCN/Search--G,
governance, experiments/results, current GROK/CLAUDE artifacts, the 2026-10-02 independent review,
and concurrent PRIMEPH1 were reconciled.

\textbf{Duplicate theorem reconciliation.}
Two concurrent LSPV-5 derivations proved the same closure at different generality/notation.
PHCLOSE1 states: if a polynomially balanced relation has graph in PH, then under \(P=NP\) a valid
witness is selectable in FP by prefix extension. PRIMEPH1 states the same theorem under the label
PRIME-PH-CLOSE1 and explicitly instantiates the project's current reconstruction prefix
\[
\exists\sigma\,\forall u\,\exists z\,\forall y\,B(x,\sigma,u,z,y).
\]
These are not distinct load-bearing theorems. The canonical project label is henceforth
\[
\boxed{\text{PRIME-PH-CLOSE1}.}
\]
PHCLOSE1 is retained as an independent proof and broader finite regression battery.

\section*{PROVED/EXACT: PRIME-PH-CLOSE1}
Assume \(P=NP\). Let \(R(x,w)\) be polynomially balanced and let its graph
\[
G_R=\{\langle x,w\rangle:R(x,w)\}
\]
belong to PH. Then on \(D_R=\{x:\exists w R(x,w)\}\) there is an FP selector \(S\) satisfying
\[
R(x,S(x)).
\]
For a prefix \(u\),
\[
E_R(x,u)\iff\exists v\,R(x,uv)
\]
remains in PH. Since \(P=NP\Rightarrow PH=P\), polynomially many prefix decisions reconstruct a
valid witness. If \(g(x,w)\) is a polynomial-time prime/CRT/provenance annotation with polynomial
output length, then \(g(x,S(x))\in FP\).

\paragraph{Canonicality corollary.}
Unique witnesses and lexicographically canonical witnesses do not escape this closure when validity
is polynomially balanced and in fixed PH.

\paragraph{Current reconstruction corollary.}
For
\[
Good(x,\sigma)
\iff
\forall u\,\exists z\,\forall y\,B(x,\sigma,u,z,y),
\qquad B\in P,
\]
we have \(Good\in\Pi_3^P\), while sigma-prefix extension lies in \(\Sigma_4^P\). Under
\(P=NP\), both collapse to P, so a valid \(\sigma\), when one exists, is selectable in FP.
This is qualitative closure only; RAMCOMP1 still controls quantitative exponent accounting.

\section*{Computational replication}
Two independent finite controls now agree.

\begin{enumerate}
\item PRIMEPH1 exhaustively tests
\[
\exists\sigma_5\forall u_3\exists z_3\forall y_2 R
\]
for unique and nonunique valid sigma sets. Five prefix queries recover the unique witness in one
case and the lexicographically first witness in the other, together with its prime annotation.
\item PHCLOSE1 tests 50 relations each with existential, \(\Pi_2\)-style, and \(\Pi_3\)-style
validity graphs on five witness bits. All 150 cases recover the lexicographically least witness in
five prefix decisions.
\end{enumerate}

Both are \textbf{COMPUTATIONAL} sanity checks only. Neither simulates \(P=NP\).

\section*{Alternation boundary}
Fixed bounded alternation remains inside PH and is closed by PRIME-PH-CLOSE1 under \(P=NP\).
Chandra--Kozen--Stockmeyer give
\[
\boxed{AP=PSPACE.}
\]
Thus genuine polynomially many alternations are the first standard syntactic zone not removed by
this closure. But a polynomial-depth challenge strategy is not one realized polynomial transcript;
it is a strategy object, potentially of exponential explicit size. Moving there is not progress
unless a model-preserving theorem forces every hypothetical polynomial static SAT decider to realize
that strategy without hiding PSPACE or exponential output in the definition.

\section*{Structural update: TURANEX1}
TURANEX1 is adopted within the structural diagnostic branch. If every \(s\)-set has
\(\epsilon(I)<r\), then
\[
\boxed{
d_r(U)\ge \binom m2-t_{s-1}(m).
}
\]
The contrapositive turns a low-support \(r\)-dimensional subcode into an explicit excess-clique
certificate. At \(r=1\), universal zero excess on \(s=\Theta(\sqrt m)\) would force
\(d_{\min}(U)=\Omega(m^{3/2})\). This remains a relativizing structural statement, not a runtime
lower bound.

\section*{Barrier and falsification ledger}
The provenance closure ladder is now:
\[
\text{semantic factoring}
\;	o\;
\text{deterministic replay}
\;	o\;
\text{NP replacement search}
\;	o\;
\text{fixed-PH replacement search}.
\]
The corresponding generic escape attempts are closed by semantic equivalence, TRACECLOSE1,
PRIME-FNP-CLOSE1, and PRIME-PH-CLOSE1 respectively.

Retain as FALSIFIED/RETIRED generic mechanisms: prime magnitude/factorization, deterministic trace
tags, canonical or unique fixed-PH prime provenance, fixed finite alternation alone, causal depth
alone, ordinary switching as universal transfer, raw local-nullity abundance, and direct
complex/real-to-binary rank transfer.

PRIME-PH-CLOSE1 relativizes and is not a nonalgebrizing hinge.

\section*{Shortest surviving proof chain}
\[
P=NP\to A\in TIME(N^K)
\to\boxed{\text{exact reconstruction-compatible SU/vMS realization}}
\to\boxed{\operatorname{rank}_{\mathbb F_2}\Phi=\omega(m)\text{ or classified exceptions}}
\]
\[
\to\boxed{\text{charged acquisition surviving REACQ1/TRACECLOSE1/PRIME-PH-CLOSE1}}
\to\boxed{\text{model-matched reconstruction}}
\to\boxed{\text{nonrelativizing/nonalgebrizing hinge}}
\to\boxed{\text{unconditional contradiction}}.
\]
Every load-bearing boxed implication remains OPEN or CONDITIONAL.

\section*{ONE SINGLE TEAM TARGET}
For one exact reconstruction-compatible SU/vMS polynomial-bit orbit, construct the source-transcribed
binary matrix \(\Phi\) and legal slice matrices \(X_s\), then prove or falsify
\[
\boxed{\operatorname{rank}_{\mathbb F_2}\Phi=\omega(m)}
\]
using intrinsically characteristic-two certificates. In the same pipeline compute odd/maximal
minors, binary dual annihilators, MATRIXCERT1 pair incidence, EXCESS/TURANEX1 generalized weights,
and COVERBREAK1 controls.

Prime provenance remains instrumentation. Permit only one further provenance falsification:
formalize the smallest source-legal polynomial-alternation SAME\_REALIZATION strategy relation and
require a model-preserving compiler from arbitrary static SAT. If that compiler fails, do not
promote the provenance branch further.

\section*{What this does not prove}
The PH closure, alternation boundary, TURANEX1, and finite experiments do not prove an unrestricted
SAT lower bound. Moving a provenance condition from PH toward PSPACE is not itself a contradiction.
P versus NP remains OPEN.

\begin{thebibliography}{9}
\bibitem{Stockmeyer76} L. J. Stockmeyer, The Polynomial-Time Hierarchy,
Theoretical Computer Science 3(1), 1976, 1--22.
\bibitem{Wrathall76} C. Wrathall, Complete Sets and the Polynomial-Time Hierarchy,
Theoretical Computer Science 3(1), 1976, 23--33.
\bibitem{CKS81} A. K. Chandra, D. C. Kozen, L. J. Stockmeyer,
Alternation, J. ACM 28(1), 1981, 114--133.
\bibitem{Krentel88} M. W. Krentel, The Complexity of Optimization Problems,
JCSS 36(3), 1988, 490--509.
\bibitem{KKMP21} R. Kleinberg, O. Korten, D. Mitropolsky, C. Papadimitriou,
Total Functions in the Polynomial Hierarchy, ITCS 2021, LIPIcs 185:44.
\bibitem{FIM25} N. Fleming, D. Imrek, C. Marciot,
Provably Total Functions in the Polynomial Hierarchy, CCC 2025, LIPIcs 339:28.
\end{thebibliography}

\end{document}
% END_OF_VERIFIED_MAIN::LSPV-5::0948
